% commentaire.tex — f168v, vv. 509–550 (conversion initiale des travaux antérieurs)
\comm{509}{τὸ γὰρ ἦτρον τῆς χύτρας ἐλάκτισεν·}{Il y a une plaisanterie \italique{contre l'attente} (παρὰ προσδοκίαν) : τῆς χύτρας est mis pour τῆς μητρός.}
\comm{515}{ἁπαξάπαντα.}{La Fontaine, \italique{La chatte métamorphosée en femme} (II, 18) :\citComm{\francais{La croit femme en tout et partout.}}}
\comm{516}{κύτταρον,}{synonyme de κῶνος, \italique{pomme de pin}, que les Grecs appelaient aussi στρόβιλος (στρόμβος dans Nicandre. Hehn., p.292). (Note de Willems.)}
\comm{518–9}{θυμούμεθα, / οὐδὲν παθοῦσαι μεῖζον ἢ δεδράκαμεν;}{On sait grâce au scholiaste que ces vers sont tirés du \italique{Télèphe} : 
\citComm{\grec{παρὰ τὰ ἐκ Τηλέφου Εὐριπίδου· `εἰ δὴ θυμούμεθα παθόντες οὐδὲν μᾶλλον ἢ δεδρακότες'.}}
(fragment 20 du \italique{Télèphe} de l'édition Jouan-van Looy aux Belles Lettres, tome VIII, \nieme{3} partie. Correspond aux fragments 711 de l'édition Kannicht et 965 de celle de Mette.)}
\comm{520–30}{}{Antistrophe du chant du chœur des vers 433-442.}
\comm{527}{ἀλλ᾽ ἅπαν γένοιτ᾽ ἂν ἤδη·}{Le lieu-commun du \latin{omnia mutantur}\nc{Cf. Tosi, \italique{Dictionnaire des sentences latines et grecques}, §26. La citation est d'Ovide (\italique{Métamorphoses}, XV, 165). Érasme (\italique{Adages}, 663) recense la formule équivalente de Térence : \latin{Omnium rerum vicissitudo est}.} avait bien-sûr cours aussi chez les Grecs. 
Brève recension : \\
Hérodote (V, 9) :\citComm{γένοιτο δ᾽ ἂν πᾶν ἐν τῷ μακρῷ χρόνῳ.}Sophocle (\italique{Philoctète}, vers 305-306) : \citComm{πολλὰ γὰρ τάδε / ἐν τῷ μακρῷ γένοιτ᾽ ἂν ἀνθρώπων χρόνῳ.}Eupolis\nc{Fragment 391(356) des \italique{Poetae Comici Graeci}, ed. Kassel \& Austin, V (1986).} :\citComm{ἦ πολλά γ' ἐν μακρῷ χρόνῳ γίγνεται μεταλλαγῇ / (τῶν) πραγμάτων· μένει δὲ χρῆμ' οὐδὲν ἐν ταὐτῷ ῥυθμῷ.}Ménandre (fragment 685(466)\nc{\italique{Poetae Comici Graeci}, VI 2 (1998). (Le fragment est extrait de Stobée I, 8, 31.)}) : 
\citComm{οὐκ ἔστ' ἄπιστον οὐδὲν ἐν θνητῷ βίῳ / οὐδ' ἂν γένοιτο.}
Et dans la pièce d'Aristophane d'antan, \italique{Lysistrata} (vers 258) : 
\citComm{Χορὸς γερόντων --- ἦ πόλλ᾽ ἄελπτ᾽ ἔνεστιν ἐν τῷ μακρῷ βίῳ φεῦ.}
Remarquons que dans tous ces extraits, le grec insiste sur la longueur de temps que prennent les changements pour s'opérer (ἐν τῷ μακρῷ χρόνῳ / ἐν θνητῷ βίῳ / ἐν τῷ μακρῷ βίῳ)... or c'est tout le contraire ici : ἤδη, \italique{aujourd'hui}, \italique{actuellement}. Ce contraste est certainement une allusion d'Aristophane aux temps mouvementés politiquement de la représentation des \italique{Thesmophories} -- où tout semble s'accélérer.}
\comm{528–30}{τὴν παροιμίαν δ᾽ ἐπαινῶ / τὴν παλαιάν· ὑπὸ λίθῳ γὰρ / παντί που χρὴ / μὴ δάκῃ ῥήτωρ ἀθρεῖν.}{Le proverbe – avec un scorpion à la place de l'orateur – était effectivement populaire puisqu'il nous est diversement parvenu dans les recueils de Proverbes. Ainsi des parémiographes Zénobie\nc{\latin{Zenobii centuria} VI, 20 (tomus I). Je cite les proverbes d'après le recensement d'E. Leutsch et F. G. Schneidewin, \italique{Corpus paroemiographorum Graecorum}, en deux tomes. Signalons qu'une nouvelle édition critique des recensions de Zénobius et Diogénianus avec traduction en italien a paru en 2019 sous le titre \italique{I proverbi greci}, dir. Emanuelle Lelli.} ou Diogénianus\nc{\latin{Diogeniani centuria} VIII, 59 (tomus I).} et d'Hésychius\nc{Hésychius (υ 717). Photius et Suidas rapportent aussi ce proverbe : il est partout !} qui rapportent la version la plus courante : 
\citComm{ὑπὸ παντὶ λίθῳ σκόρπιος εὕδει: ἐπὶ τῶν κακοήθων λέγεται.}Macaire\nc{\latin{Macarii centuria} VIII, 69 (tomus II).} en précise le sens : \grec{ἐπὶ πολυπληθίᾳ τῶν πανταχοῦ ὄντων πονηρῶν ἀνθρώπων}, « rapport au grand nombre d'hommes qui sont méchants sous tout rapport ».
Le proverbe est attesté aussi par Praxilla\nc{Le scholiaste attribue même le proverbe à Praxilla. Athénée (XV, 695d) donne un exemple de scolie (σκόλιον), sorte de chanson à boire, écrite sur ce proverbe ; comme Praxilla excella dans ce genre, il n'est pas impossible que la scolie fût d'elle, et que le scholiaste l'ait en tête lorsqu'il attribue un peu hâtivement la maternité du proverbe à Praxilla. Cette scolie correspond au fragment 903 des \italique{Poetae Melici Graeci}.} et Sophocle\nc{Fragment 37 de l'édition Radt.}.\\
Ici, la \italique{pierre} (λίθος) dont parle la coryphée est le piédestal (βῆμα) de la Pnyx sur lequel les orateurs -- Mnésiloque, en l'occurrence -- parlaient. (Aristophane use du mot dans le même sens aux vers 683 des \italique{Acharniens} ; 680 de la \italique{Paix} et 87 de l'\italique{Assemblée des femmes}.)
Hésychius (λ 999) :\citComm{λίθος· βῶλος. βωμός. καὶ βάσις. τὸ ἐν τῇ Ἀθηναίων ἐκκλησίᾳ βῆμα.}Enfin, la substitution au scorpion de l'orateur n'est pas gratuite : on parlait couramment de la \italique{morsure de la langue}\nc{Voir les sens figurés du verbe δάκνειν.}.
Ainsi de la \italique{morsure du sycophante}, δῆγμα συκοφάντου, qui paraît proverbiale\nc{Voir le \italique{Ploutos}, vers 885, et le commentaire du scholiaste \italique{ad loc.}}. Diogène le cynique eut cette réplique à l'ironie \italique{mordante} (rapportée par son homonyme laërtien) :
\citComm{ἐρωτηθεὶς τί τῶν θηρίων κάκιστα δάκνει ἔφη: τῶν μὲν ἀγρίων συκοφάντης, τῶν δὲ ἡμέρων κόλαξ.\\
Interrogé sur la bête dont la morsure est la plus mauvaise, [il répondit] que des bêtes sauvages c'était le sycophante, et des animaux domestiques, c'était le flatteur.}}
\comm{531–73}{}{Tout ce passage est composé en tétramètres iambiques catalectiques ; mètre du conflit et des débats animés. Cf. aussi les \italique{Cavaliers} v. 335-366 ; \italique{Lysistrata} v. 350-381, 467-475.}
\comm{531–2}{ἀλλ᾽ οὐ γάρ ἐστι τῶν ἀναισχύντων φύσει γυναικῶν / οὐδὲν κάκιον εἰς ἅπαντα πλὴν ἄρ᾽ εἰ γυναῖκες.}{Oct. Feuillet, \italique{Dalila}, partie IV :
\citComm{\francais{Rien de plus semblable sous le soleil qu'une femme – et une femme.}}C'est la traduction presque textuelle de Ménandre (\italique{Sentences monostiques}, Γνῶμαι μονόστιχοι 109) :\citComm{Γυνὴ γυναικὸς πώποτ᾽ οὐδὲν διαφέρει.}(Note de Willems ; pour le texte, voir le commentaire critique.)\\
Il serait piquant qu'Aristophane mette dans la bouche de la coryphée une citation d'Euripide lui-même pour dénigrer le comportement de certaines femmes ; le scholiaste n'en a malheureusement pas soufflé mot... on dispose néanmoins de citations d'Euripide très approchantes (fr. 494.27 Kannicht $=$ fr. 14.27 \italique{Mélanippe} Jouan-van Looy) : \citComm{τῆς μὲν κακῆς κάκι[ο]ν οὐδὲν γίγνεται / γυναικός.}}
\comm{533}{μὰ τὴν Ἄγραυλον.}{Scholie d'un érudit :\citComm{κατὰ τῆς ᾿Αγραύλου ὤμνυον, κατὰ δὲ τῆς Πανδρόσου σπανιώτερον· κατὰ δὲ τῆς Ἕρσης οὐχ εὑρήκαμεν.\\
\francais{L'on jurait sur Agraule, plus rarement sur Pandrosos ; nous n'avons pas trouvé d'exemple sur Ersès.}} Aglaure (pour le texte, voir le commentaire critique), Pandrose et Hersé étaient les trois filles de Kécrops, mythique fondateur et premier roi d'Athènes.}
\comm{549–50}{Μίαν γὰρ οὐκ ἂν εἴποις / τῶν νῦν γυναικῶν Πηνελόπην, Φαίδρας ἁπαξαπάσας.}{Le comique Euboulos paraît s'être inspiré de ce passage :\citComm{ὦ Ζεῦ πολυτίμητ', εἶτ' ἐγὼ κακῶς ποτε\\
ἐρῶ γυναῖκας ; νὴ Δί' ἀπολοίμην ἄρα\\
πάντων ἄριστον κτημάτων. Εἰ δ' ἐγένετο\\
κακὴ γυνὴ Μήδεια, Πηνελόπη δέ <γε>\\
μέγα πρᾶγμ'. ᾿Ερεῖ τις ὡς Κλυταιμνήστρα κακή·\\
Ἄλκηστιν ἀντέθηκα χρηστήν. ᾿Αλλ' ἴσως\\
Φαίδραν ἐρεῖ κακῶς τις · ἀλλὰ νὴ Δία\\
χρηστή τις ἦν μέντοι - τίς ; οἴμοι δείλαιος,\\
ταχέως γέ μ' αἱ χρησταὶ γυναῖκες ἐπέλιπον,\\
τῶν δ' αὖ πονηρῶν ἔτι λέγειν πολλὰς ἔχω.\\
\francais{Zeus très honoré, moi, dire du mal des femmes ? Ah! plutôt mourir ! Des femmes, le plus précieux de tous les biens ! Car si Médée s'est montrée méchante femme, Pénélope, elle, fut un vrai trésor. On dira que Clytemnestre fut mauvaise ? Je mets en regard la vertueuse Alceste. Mais peut-être médira-t-on de Phèdre ? Eh, par Zeus, vertueuse sans contredit fut.... qui donc ? Ah! malheureux ! les femmes vertueuses m'ont de suite fait faute, quand des mauvaises j'en ai encore un tas à citer.}}(Meineke, \italique{Fragm. Com. Gr}., t. III, p. 260. [$=$ fragment 115 (116. 117) de l'édition Kassel-Austin, \italique{Poet. Com. Gr.}]) (Note de Willems.)}
\comm{550}{Πηνελόπην, Φαίδρας.}{Si Phèdre est souvent citée par Aristophane comme (anti-)modèle d'impudicité (voir le vers 153 ; les \italique{Grenouilles} v. 1043 ; fr. 469(453) K.-A.(Kock)), Mélanippe apparaît seulement ici dans les œuvres conservées d'Aristophane.}
